% ===================================================================
% Draft subsection for the paper: hardware validation on the VC709.
% Every \HW{...} is a placeholder. Replace it ONLY with a value from
% results/<timestamp>/ (run.log, CSVs) or reports/clk200/ (Vivado).
% Then \input the generated results/<timestamp>/tables.tex.
% ===================================================================
\providecommand{\HW}[1]{\textcolor{red}{[#1]}}

\subsection{Hardware Validation on a Virtex-7 FPGA with EEG Data}
\label{sec:hw}

The three architectures (A: direct form, B: symmetric-folded, C: folded and
pipelined) were implemented together on a Xilinx VC709 board
(XC7VX690T-2FFG1761) at \HW{200}~MHz, with the RTL unchanged from the
simulation study. A hardware-in-the-loop harness streams an EEG record from
on-chip block RAM through all three filters in parallel, one sample per clock,
stores every output in block RAM, and returns it to the host over USB-UART.
On-chip counters record the output count, the first-output latency and the
total clock count of each filter. The routed design meets timing with a
worst negative slack of \HW{WNS}~ns and uses no DSP48 slices
(\HW{dsp\_total}~=~0).

The input is the seven-epoch synthetic EEG of Agarwal et
al.~\cite{agarwal2017} (140{,}007 samples), clean and with white Gaussian
noise at input SNRs from 30 to $-5$~dB. Each record is peak-normalised to
$[-1,1]$ and quantised to 8-bit two's complement, as in the normalisation
step of AutoFIR~\cite{liu2026autofir}. Table~\ref{tab:hw_eeg} lists the
signal metrics computed from the hardware output. For every record, every
output of all three filters matched the bit-exact software model
(\HW{0} mismatches in \HW{7,560,810} compared outputs, two input rates). The
measured SNR improvement saturates at \HW{5.8}~dB, in agreement with the
analytic white-noise gain of the quantised taps,
$10\log_{10}\!\bigl[(\sum h)^2/\sum h^2\bigr] = 5.90$~dB. The fidelity with
respect to the floating-point filter is limited to about \HW{36}~dB by
coefficient rounding ($\sum h_q = 130$ instead of 128), not by the MRPM
datapath, which is exact.

Table~\ref{tab:hw_timing} gives the measured timing. With one sample per
clock the filters process the full record in \HW{140,016}/\HW{140,020}
clocks (\HW{700}~$\mu$s at 200~MHz), i.e.\ a sustained throughput of one
output per clock. The measured input-to-output latency is \HW{2} clocks for
A and B and \HW{5} clocks for C; the difference of three clocks equals the
pipeline registers added in C. The die temperature during the runs, read
from the on-chip XADC, was \HW{T}~\textdegree C at
VCCINT~=~\HW{V}~V.

% \input{results/<timestamp>/tables.tex}
